\section*{Conclusion}
This chapter began with a puzzle: % AF: non, on ne pouvait pas avoir en tête le puzzle avoir de connaître les résultats. Cf. le plan de la conclusion que j'ai proposé
internationally redistributive policies remain politically marginal even though citizens are not systematically opposed to transferring resources across borders, coordinating policy effort internationally, or delegating authority beyond the nation-state. The literature reviewed here shifts the explanation away from a simple presumption that distributive preferences are nationally bounded. Moral universalism provides a meaningful, although unevenly distributed, basis for extending distributive concern beyond compatriots \citep{enke_moral_2023}. This disposition does not translate mechanically into support for any particular policy, but it makes cross-border redistribution politically conceivable. When attention moves from abstract values to concrete instruments, the cumulative evidence supports a qualified conclusion: foreign aid, international taxation, debt relief, climate finance, and institutional cooperation can attract majority or plurality support when their objectives, incidence, and implementation are clearly specified \citep{fabre_majority_2025,fabre_public_2026}.

The successive strands of evidence point to a common interpretation. Citizens appear willing to endorse international redistribution not as an unconditional transfer of resources, but as a rule-governed collective undertaking. Support is generally stronger when beneficiaries are identifiable, resources are protected against diversion, burdens fall on those with greater capacity to contribute, and international participation is reciprocal and credible \citep{fabre_majority_2025,fabre_public_2026,gampfer_obtaining_2014}. The same logic extends from fiscal instruments to political institutions. Unspecified transfers of authority beyond the nation-state generate resistance, whereas democratic and functionally limited forms of global governance attract substantially greater acceptance \citep{ghassim_who_2024}. Similarly, making private costs explicit reduces support but does not eliminate it when policies combine environmental effectiveness, progressive incidence, and visible poverty reduction \citep{bechtel_mass_2013,fabre_public_2026}. Citizens therefore appear to evaluate not international redistribution in isolation, but the institutional package through which it is financed, allocated, implemented, enforced, and justified.

Returning to the contrast introduced at the outset, the moral and geographical distance separating Clara and Moussa does not constitute an absolute barrier to redistribution. The decisive question is whether concern for Moussa can be translated into a policy that Clara perceives as fairly financed, effectively implemented, and supported by legitimate and credible institutions.

The theoretical implication is that preferences over international redistribution cannot be adequately represented along a single cosmopolitan--nationalist dimension. They are multidimensional and conditional on interactions among moral concern, material incidence, expected effectiveness, institutional legitimacy, and beliefs about collective participation. Fairness operates vertically through the distribution of costs across income groups and horizontally through the allocation of responsibilities across countries. Effectiveness matters because citizens distinguish transfers that plausibly reach their intended beneficiaries from those exposed to corruption, free-riding, or weak enforcement. Institutional credibility connects these dimensions: democratic procedures, transparent revenue allocation, and identifiable participation can help translate a general willingness to assist distant beneficiaries into acceptance of a specific policy. This framework also clarifies why self-interest matters without being determinative. Higher private costs and competition with domestic spending constrain support, yet they coexist with substantial willingness to contribute when policies are perceived as equitable, effective, and institutionally credible.

Important tensions nevertheless remain. Much of the evidence relies on stated preferences elicited after respondents receive information and evaluate policy designs that are more explicit than those ordinarily encountered in electoral debate. Majority acceptance in surveys therefore need not imply intense preferences, sustained attention, or a willingness to sanction parties that neglect these issues. The findings also reveal substantial heterogeneity across countries, ideological groups, and income positions, while information treatments produce context-dependent effects. Correcting misperceptions increases donations or support for foreign aid in some settings, but information about citizens' global income rank does not consistently alter redistributive preferences \citep{nair_misperceptions_2018,fehr_your_2019}. Moreover, willingness to contribute may remain politically latent when individuals underestimate its prevalence, as the gap between first- and second-order beliefs in climate cooperation illustrates \citep{andre_globally_2024,mildenberger_beliefs_2017}. These limitations do not overturn the chapter's central conclusion; rather, they identify mechanisms through which favorable attitudes may fail to become organized political demand.

The central unresolved question is therefore less whether a constituency for international redistribution exists than under what conditions its latent and conditional acceptance becomes politically consequential. Future research should examine how issue salience is created, how citizens update beliefs about others' preferences, how parties and representatives infer electoral demand, and which institutional features sustain support outside survey settings. It should also distinguish more systematically among policy approval, willingness to pay, budgetary prioritization, political participation, and vote choice, each of which captures a different stage in the translation of moral concern into public policy. This agenda redirects attention from the measurement of isolated attitudes toward the political processes connecting public preferences, institutional design, and policy supply. Explaining this translation---and the conditions under which it succeeds or fails---is the next step toward understanding why internationally redistributive policies remain scarce despite a broader social basis of acceptance than prevailing political debates suggest.


